% II. Анализ на възможните решения и обосновка на избора.
% По едно решение на всяко отворено място: език, кодиране на термините, набор
% от индекси, носител на индексите, вид на синтактичния анализатор, стратегия
% на планировчика.
\section{Анализ на възможните решения и обосновка на избора}
\label{sec:alternatives}

\subsection{Целеви език и среда за изграждане}

Разгледани са три възможности: C++20, Rust и Java. Java отпада заради изискването
хранилището да бъде вградено в приложение на друг език без отделна среда за
изпълнение. Rust покрива изискването и предлага по-строги гаранции за паметта, но
свързването му към съществуващ код на C++ минава през допълнителен слой. Избран е
C++20, защото целевото приложение на такава библиотека най-често е на C++, защото
езикът дава пряк контрол върху разположението на данните в паметта, което е
определящо за индексите, и защото авторът има опит в него. Изграждането е с CMake,
тъй като това е фактическият стандарт за преносими проекти на C++.

\subsection{Кодиране на термините}

Триплетът може да съхранява самите низове на термините или числови идентификатори,
получени от речник. Прякото съхранение на низове е по-просто, но повтаря един и същ
IRI толкова пъти, колкото триплета го използват, а сравненията при съединяване стават
низови. Речниковото кодиране изисква два допълнителни справочника, от термин към
идентификатор и обратно, но свежда триплета до три числа с фиксиран размер. Това
прави индексите плътни, сравненията евтини, а съединенията сводими до сливане на
сортирани редици. Избрано е речниково кодиране, следвайки~\cite{neumann2010rdf3x}.

\subsection{Набор от индекси}

Пълният набор от шест пермутации обслужва всяка комбинация от свързани позиции с
подходящо подреден индекс~\cite{weiss2008hexastore}, но умножава обема шест пъти.
Един единствен индекс по подлог е компактен, но при шаблон със свободен подлог води
до пълно обхождане. Избрани са три пермутации, а именно подлог сказуемо допълнение,
сказуемо допълнение подлог и допълнение подлог сказуемо. Всяка от осемте комбинации
от свързани и свободни позиции се обслужва от префикс на поне един от трите индекса,
което дава покритие без шесткратното повторение. Обосновката е обобщена в
\tabref{tab:alternatives}.

\subsection{Носител на индексите}

Възможностите са структури изцяло в паметта, вграден ключ и стойност склад като
вариант на дърво с логаритмично сливане, и собствени файлове, отразени в паметта.
Структурите в паметта са най-бързи, но задължават цялият набор да се събира в
оперативната памет. Вграден външен склад решава това, но внася зависимост, чийто
модел на съхранение не съвпада с подредените редици, които съединяването изисква.
Избрани са собствени файлове, отразени в паметта: подредбата на записите остава под
контрола на проекта, операционната система поема кеширането, а отварянето на
съществуващо хранилище не изисква повторно изграждане на индексите.

\subsection{Вид на синтактичните анализатори}

Граматиките на Turtle и на SPARQL~1.1 са публикувани в спецификациите и могат да
бъдат подадени на генератор на анализатори. Генераторът пести писане и намалява риска
от разминаване с граматиката, но добавя зависимост, стъпка в изграждането и
генериран изходен текст, чиито съобщения за грешка са трудни за привеждане към
полезен вид. Двете граматики са проектирани така, че да се разпознават с ограничен
предварителен поглед, което ги прави удобни за анализатор с рекурсивно спускане,
писан на ръка. Избран е вторият вариант, като рискът от разминаване с граматиката се
покрива от тестовите набори на W3C~\cite{w3c_sparql11_testsuite}.

\begin{table}[H]
\caption{Обобщение на разгледаните алтернативи и на направения избор}
\label{tab:alternatives}
\centering
\fontsize{10}{12}\selectfont
\begin{tabular}{@{}L{3.0cm}L{4.2cm}L{3.0cm}L{3.2cm}@{}}
\toprule
\textbf{Решение} & \textbf{Разгледани варианти} & \textbf{Избор} & \textbf{Водещ довод} \\
\midrule
Език & C++20, Rust, Java & C++20 & вграждане без отделна среда за изпълнение \\
Термини & низове в триплета, речниково кодиране & речниково кодиране & фиксиран размер, евтини сравнения \\
Индекси & един, три, шест пермутации & три пермутации & покритие на всички шаблони без шесткратен обем \\
Носител & памет, външен склад, отразени файлове & отразени файлове & контрол върху подредбата, без нова зависимост \\
Анализатори & генератор, рекурсивно спускане & рекурсивно спускане & по-добри съобщения, без стъпка в изграждането \\
\bottomrule
\end{tabular}
\end{table}

\subsection{Стратегия на планировчика}

Редът на съединенията при базов графов шаблон определя размера на междинните
резултати и с това времето за изпълнение. Разгледани са два подхода: фиксирано
правило, което подрежда шаблоните по брой свързани позиции, и избор по оценка,
основан на статистики от самите индекси. Първият е тривиален за реализация и няма
цена по време на планиране, но е чувствителен към реда на изписване на заявката.
Вторият изисква поддържане на статистики, но остава устойчив. Избран е вторият, като
статистиките се извличат от дължините на префиксите в индексите.
\TODO{решение дали статистиките се поддържат постъпково при вмъкване или се
преизчисляват при отваряне на хранилището}
